%=============================================================================!
% 8.6  MANY-BODY POTENTIALS                                                   !
%=============================================================================!
\section{Many-body potentials for metals and carbon}
\label{sec:pe-manybody}

In a pair potential every bond has the same strength whatever else its
atoms are bonded to. That fails for two materials used for the negative
electrodes of lithium batteries, lithium metal and graphite. In
both, the strength of a bond depends on how many other bonds its atoms
have. This section builds a model with that property, the simplest
many-body potential, and describes the forms used for metals and carbon.

\subsection*{Bonds that weaken as they multiply}

In a metal the outer electrons are shared by all the atoms, and an atom
with more neighbours spreads its share of them over more bonds, each
weaker. A model that captures this writes the energy of atom $i$ with a
square root,
\begin{equation}
   U = \sum_i\Bigl[\sum_{j\ne i}A\,\mathrm{e}^{-\eta(r_{ij}/r_0 - 1)}
   - \sqrt{\rho_i}\,\Bigr] ,
   \qquad
   \rho_i = \sum_{j\ne i}\xi^2\,\mathrm{e}^{-2\zeta(r_{ij}/r_0 - 1)} ,
   \label{eq:pe-second-moment}
\end{equation}
with a pair repulsion of strength $A$ and range set by $\eta$, and a bond
energy $-\sqrt{\rho_i}$ built from a sum $\rho_i$ over the neighbours,
each contributing $\xi^2$ at the distance $r_0$ and less farther away.
Finnis and Sinclair proposed this square root for metals such as iron
and tungsten\cite{FinnisSinclair1984}, and the exponential forms used here are
those of Cleri and Rosato\cite{CleriRosato1993}; it is called the
second-moment model.

For an atom with $z$ neighbours, all at $r_0$, $\rho_i = z\xi^2$ and its
bond energy is $-\xi\sqrt{z}$: it grows only as the square root of the
number of bonds, so each bond is worth $-\xi/\sqrt{z}$, weaker the more
there are. In a pair potential the bond energy is $z$ times that of one
bond. In a close-packed crystal, its atoms stacked like oranges on a market
stall, each atom touches 12 others; in the body-centred cubic structure of
lithium each sits at the centre of a cube of 8; and in a simple cubic
crystal each has 6, along the three axes. \Cref{fig:pe-manybody} compares
the two models, scaled to agree for close-packed metals, $z = 12$. A single bond of the second-moment model
is worth $1/\sqrt{12} = 0.289$ of a lone bond ($z = 1$) when the atom has
12 neighbours,
and $1/\sqrt{8} = 0.354$ in body-centred cubic lithium, with 8. One
consequence is that an atom at a surface, with fewer neighbours, holds
those it has more tightly, and \cref{ex:pe-vacancy} shows that, counting
only the bond energy, making a vacancy, by moving an atom from inside a
close-packed crystal to its surface, costs about half the energy that
holds an atom in the crystal, where a pair potential makes the two
equal.

\begin{figure}[tbp]
\centering
\includegraphics{ch08_potentials/manybody}
\caption{The bond energy of an atom with $z$ neighbours, as a fraction of
its value at $z = 12$: proportional to $z$ for a pair potential
(dashed), and to $\sqrt{z}$ for the second-moment model (teal), at a pair
of atoms (1), a chain (2), a graphite layer (3), diamond (4), simple
cubic (6), body-centred cubic lithium (8) and close-packed metals (12).}
\label{fig:pe-manybody}
\end{figure}

The forces follow from the chain rule. The square root couples every
neighbour of $i$: moving atom $k$ changes $\rho_i$ through the term
$r_{ik}$, and $\sqrt{\rho_i}$ changes by $\frac12\rho_i^{-1/2}$ times that.
The force on an atom therefore depends on the neighbours of its
neighbours, and is no longer a sum of forces along pairs with fixed
strengths. The forces still add to zero and exert no total torque, because $U$
depends only
on distances and so does not change when all the atoms are moved or
turned together (\cref{sec:ha-noether}); the function
\texttt{second\_moment} of \texttt{mdlab} computes it, and its test
checks it against central differences.

The embedded-atom method generalises \cref{eq:pe-second-moment}: the
energy of atom $i$ is $\mathcal{F}(\bar\rho_i) + \frac12\sum_j\varphi
(r_{ij})$, with
$\bar\rho_i$ the density of the electrons, their charge per unit volume,
that its neighbours produce where it sits and $\mathcal{F}$, the embedding energy, a function fitted to measurements or
to calculations\cite{DawBaskes1984}. The second-moment model is the case
$\mathcal{F}(\bar\rho) = -\sqrt{\bar\rho}$, with $\bar\rho_i = \rho_i$ and
$\varphi(r) = 2A\,\mathrm{e}^{-\eta(r/r_0 - 1)}$, the 2 making up for the
$\frac12$.

\subsection*{Carbon and the order of a bond}

Carbon forms directional bonds, three to a layer in graphite and four in
diamond, and its bonds too weaken as they multiply. Tersoff's potential
writes the energy as a sum over bonds, $\frac12\sum_i\sum_{j\ne i}f_{\mathrm
c}(r_{ij})\bigl[f_{\mathrm R}(r_{ij}) - b_{ij}f_{\mathrm A}(r_{ij})
\bigr]$, with a repulsion $f_{\mathrm R}$, an attraction $f_{\mathrm A}$,
a smooth cutoff $f_{\mathrm c}$, and a bond order $b_{ij}$ that falls as
atom $i$ gains other neighbours, depending also on the angles between its
bonds\cite{Tersoff1988,Tersoff1988carbon}. The reactive bond-order potential of Brenner and
co-workers extends this to hydrocarbons and describes the strong bonds
within a graphite layer\cite{Brenner2002}. It has no interaction between
the layers, which are held together by the weak attraction between
neutral atoms; a version adds a Lennard-Jones term for
it\cite{Stuart2000}. Lithium between the layers of graphite needs, in
addition, lithium-carbon and lithium-lithium interactions, and models
fitted for this are among the machine-learned potentials planned for the
continuation of Part~V.

\begin{keyidea}
In metals and in carbon the strength of a bond falls as an atom's bonds
multiply. The second-moment model, with bond energy $-\sqrt{\rho_i}$,
makes the bond energy grow as $\sqrt{z}$; the embedded-atom method and
bond-order potentials refine the same idea. Their forces couple the
neighbours of neighbours but still add to zero and exert no total
torque.
\end{keyidea}

\notebook{08}{the second-moment model: bonds against neighbours}

\begin{selfcheck}
In the second-moment model, by what factor is one bond of an atom in a
chain ($z = 2$) stronger than one bond in close-packed metal?
\end{selfcheck}
